%% ============================================================
%% Appendix C: Eigenvalue Sensitivity and SVD
%% ============================================================

\chapter{Eigenvalue Sensitivity and Singular Value Decomposition}
\label{app:eigenvalue}

This appendix covers two topics moved from Chapter~5 to keep the
main text focused: sensitivity of eigenvalues and eigenvectors
(important for stability analysis), and the singular value
decomposition (important for understanding ill-conditioning).

%%------------------------------------------------------------
\section{Sensitivity of Eigenvalues}\index{eigenvalue sensitivity}
\label{appsec:eigensens}
%%------------------------------------------------------------

Let $A(\varepsilon) = A_0 + \varepsilon B$ be a matrix depending
on a scalar parameter $\varepsilon$. The eigenvalue $\lambda_k(\varepsilon)$
and eigenvector $\vec{x}_k(\varepsilon)$ satisfy
$A(\varepsilon)\vec{x}_k = \lambda_k\vec{x}_k$.

\subsection{Symmetric Matrices}

For symmetric $A_0$ with distinct eigenvalues, differentiating
$A\vec{x}_k = \lambda_k\vec{x}_k$ at $\varepsilon = 0$:
\begin{equation}
  \dd{\lambda_k}{\varepsilon}\bigg|_{\varepsilon=0}
  = \vec{x}_k^T B\vec{x}_k.
  \label{eq:eigsens_sym}
\end{equation}
The eigenvalue sensitivity is a quadratic form in the perturbation
matrix and the normalized eigenvector.

\subsection{Non-Symmetric Matrices}

For non-symmetric $A_0$, both left eigenvectors $\vec{y}_k$
(satisfying $\vec{y}_k^T A = \lambda_k\vec{y}_k^T$) and right
eigenvectors $\vec{x}_k$ are needed:
\begin{equation}
  \dd{\lambda_k}{\varepsilon}\bigg|_{\varepsilon=0}
  = \frac{\vec{y}_k^T B \vec{x}_k}{\vec{y}_k^T \vec{x}_k}.
  \label{eq:eigsens_nonsym}
\end{equation}
The formula is well-defined when $\lambda_k$ is a simple (distinct)
eigenvalue, because in that case $\vec{y}_k^T\vec{x}_k \neq 0$.

\begin{proposition}
\label{prop:biorthogonality}
If $A_0$ has distinct eigenvalues $\lambda_k \neq \lambda_j$, then
the corresponding left and right eigenvectors satisfy
$\vec{y}_k^T\vec{x}_j = 0$ for $k \neq j$
(\emph{biorthogonality}) and $\vec{y}_k^T\vec{x}_k \neq 0$.
\end{proposition}

\begin{proof}
Let $A\vec{x}_j = \lambda_j\vec{x}_j$ and $\vec{y}_k^T A = \lambda_k\vec{y}_k^T$.
Multiply the first equation on the left by $\vec{y}_k^T$ and the
second on the right by $\vec{x}_j$:
\[
  \vec{y}_k^T A \vec{x}_j = \lambda_j(\vec{y}_k^T\vec{x}_j)
  \qquad\text{and}\qquad
  \vec{y}_k^T A \vec{x}_j = \lambda_k(\vec{y}_k^T\vec{x}_j).
\]
Subtracting: $(\lambda_k - \lambda_j)(\vec{y}_k^T\vec{x}_j) = 0$.
Since $\lambda_k \neq \lambda_j$, we conclude $\vec{y}_k^T\vec{x}_j = 0$.

To see that $\vec{y}_k^T\vec{x}_k \neq 0$: if it were zero, then
$\vec{x}_k$ would be orthogonal to every $\vec{y}_j$ (via biorthogonality
for $j \neq k$ and the assumed zero for $j = k$). But the left
eigenvectors $\{\vec{y}_j\}$ span $\mathbb{R}^n$ when eigenvalues are
distinct~\cite[Theorem~1.4.5]{horn2012matrix}, so no nonzero vector
can be orthogonal to all of them. This contradiction shows
$\vec{y}_k^T\vec{x}_k \neq 0$.
\end{proof}

When $\lambda_k$ is \emph{not} distinct (a repeated eigenvalue),
$\vec{y}_k^T\vec{x}_k$ can be zero, and the formula~\eqref{eq:eigsens_nonsym}
breaks down — this is the mathematical signature of a defective matrix
and, in applications, a bifurcation point.

When $\vec{y}_k^T\vec{x}_k$ is small but nonzero, the eigenvalue is highly
sensitive to perturbations. This occurs near a bifurcation where
two eigenvalues nearly coincide.

For the SIR Jacobian at the endemic equilibrium near $R_0 = 1$,
the left and right eigenvectors corresponding to the near-zero
eigenvalue become nearly orthogonal as $R_0 \to 1$, explaining
the condition-number spike observed in Chapter~8.

%%------------------------------------------------------------
\section{Singular Value Decomposition}
\label{appsec:SVD}
%%------------------------------------------------------------

Every matrix $A \in \R^{m\times n}$ has the \textbf{singular value
decomposition} (SVD):
\begin{equation}
  A = U\Sigma V^T,
  \label{eq:SVD}
\end{equation}
where $U \in \R^{m\times m}$ and $V \in \R^{n\times n}$ are orthogonal
and $\Sigma \in \R^{m\times n}$ is diagonal with non-negative entries
$\sigma_1 \geq \sigma_2 \geq \cdots \geq \sigma_r > 0 = \sigma_{r+1} = \cdots$
(the singular values).

\subsection{SVD and the Condition Number}\index{condition number!SVD definition}

The condition number from Chapter~8 is
$\kappa_2(A) = \sigma_1/\sigma_r$.
The SVD makes clear why: the smallest singular value $\sigma_r$
measures how close $A$ is to being rank-deficient. When $\sigma_r \approx 0$,
$A$ is nearly singular and tiny perturbations to the right-hand side
$\vec{b}$ can produce large changes in the solution $\vec{u} = A^{-1}\vec{b}$.

\subsection{SVD and SA}

The sensitivity Jacobian $\mathbf{S}$ from Chapter~3 can be analyzed
using its SVD. The singular vectors of $\mathbf{S}$ identify the
combinations of parameters that most strongly affect the model output
(left singular vectors) and the combinations of outputs most strongly
affected (right singular vectors).

This connects to Principal Component Analysis (PCA):
the right singular vectors of the design matrix $X$ in a regression
are the principal components. If the SA study uses LHS sampling
(Chapter~9), PCA of the sample matrix reveals parameter combinations
that drive most of the variation in model output, providing global
SA insights that are complementary to Sobol indices.

\subsection{SVD and the Pseudoinverse}\index{pseudoinverse}

When $A$ is not square or is rank-deficient, the Moore-Penrose
pseudoinverse $A^+ = V\Sigma^+ U^T$ provides the minimum-norm
least-squares solution: $\vec{u} = A^+\vec{b}$.
The pseudoinverse appears in parameter estimation from
sensitivity data and in regularized inversion methods.

%%------------------------------------------------------------
\section{Exercises}
\label{appsec:C_exercises}
%%------------------------------------------------------------

\begin{selfcheck}
The sensitivity Jacobian $\mathbf{S}$ for the SIR model at $T = 60$~days
(from the Chapter~3 MATLAB laboratory) is a $3\times4$ matrix.
(a) Explain what the SVD of $\mathbf{S}$ reveals about the SA problem
that the individual sensitivity indices do not.
(b) If the largest singular value of $\mathbf{S}$ is $\sigma_1 = 2.3$
and the smallest is $\sigma_4 = 0.02$, what is the condition number
$\kappa_2(\mathbf{S})$? Interpret this number in plain language for
the SA context.
(c) The right singular vector $\vec{v}_1$ corresponding to $\sigma_1$
identifies the most ``influential'' linear combination of \POIs.
If $\vec{v}_1 \approx (0.7, 0.7, 0.1, 0.0)^T$ for \POIs
$(k, \beta, \tau, \mu)$, what does this tell you about the
SA of the SIR model?
\end{selfcheck}
\begin{scAnswer}
(a) The SVD reveals which \emph{combinations} of parameters drive the
largest output changes. Individual SIs only measure one-at-a-time effects;
the SVD captures the directions in parameter space that most amplify
perturbations to the model output.
(b) $\kappa_2 = 2.3/0.02 = 115$. This means that a perturbation in
the direction $\vec{v}_4$ (least influential parameter combination)
has only $1/115$ the output effect of a perturbation in direction $\vec{v}_1$.
The SA result is dominated by a small subspace of parameter space.
(c) The most influential combination weights $k$ and $\beta$ almost equally
($0.7\approx0.7$) with little contribution from $\tau$ and essentially none
from $\mu$. This confirms the structural correlation identified in Chapter~3:
$k$ and $\beta$ appear only as the product $k\beta$ in the SIR model,
so they always have equal influence on any output.
\end{scAnswer}

\begin{selfcheck}
For a symmetric $2\times2$ matrix $A = \begin{pmatrix}a&b\\b&d\end{pmatrix}$
with $a>d>0$ and $b$ small, the eigenvalue sensitivity formula~\eqref{eq:eigsens_sym}
gives $d\lambda_1/d\varepsilon = \vec{x}_1^T B \vec{x}_1$
where $B$ is the perturbation matrix.
(a) Find the eigenvectors of $A$ for $b=0$.
(b) For the perturbation $B = \begin{pmatrix}0&1\\1&0\end{pmatrix}$
(an off-diagonal perturbation), compute $d\lambda_1/d\varepsilon$ at $b=0$.
(c) Near a transcritical bifurcation in the SIR model, two eigenvalues
of the Jacobian approach each other. Why does formula~\eqref{eq:eigsens_nonsym}
become unreliable in this case?
\end{selfcheck}
\begin{scAnswer}
(a) For $b=0$: eigenvectors are $\vec{e}_1 = (1,0)^T$ ($\lambda_1 = a$)
and $\vec{e}_2 = (0,1)^T$ ($\lambda_2 = d$).
(b) $d\lambda_1/d\varepsilon = \vec{x}_1^T B \vec{x}_1 = (1,0)\begin{pmatrix}0&1\\1&0\end{pmatrix}(1,0)^T = 0$.
The leading eigenvalue has zero first-order sensitivity to this off-diagonal perturbation
(by symmetry; it appears at second order).
(c) When two eigenvalues nearly coincide ($\lambda_k \approx \lambda_j$),
the eigenvectors $\vec{x}_k$ and $\vec{x}_j$ become nearly parallel,
making $\vec{y}_k^T\vec{x}_k \approx 0$ in the denominator of~\eqref{eq:eigsens_nonsym}.
The sensitivity diverges, signaling that eigenvalue perturbation theory breaks down,
the same mathematical mechanism that causes SA indices to diverge near a bifurcation
(Chapter~8).
\end{scAnswer}
